% Bien dich bang XeLaTeX
\documentclass[a4paper]{article}

\usepackage{fontspec}
\setmainfont{Times New Roman}

\usepackage[left=3.0cm,right=2.4cm,top=2.8cm,bottom=3.0cm]{geometry}
\usepackage{setspace}
\onehalfspacing

\usepackage{amsmath}
\usepackage{amssymb}
\usepackage{booktabs}
\usepackage{longtable}
\usepackage{tabularx}
\usepackage{array}
\usepackage{enumitem}
\usepackage{titlesec}
\usepackage{indentfirst}
\usepackage{caption}
\usepackage{fancyhdr}
\usepackage[dvipsnames]{xcolor}
\usepackage{xurl}
\usepackage[unicode,colorlinks=true,linkcolor=NavyBlue,citecolor=NavyBlue,urlcolor=Blue]{hyperref}

\makeatletter
\renewcommand\normalsize{%
    \@setfontsize\normalsize{13pt}{19.5pt}%
    \abovedisplayskip 13pt plus 3pt minus 7pt%
    \abovedisplayshortskip \z@ plus 3pt%
    \belowdisplayshortskip 7pt plus 3.5pt minus 4pt%
    \belowdisplayskip \abovedisplayskip
    \let\@listi\@listI
}
\makeatother
\normalsize

\setlength{\parindent}{1.27cm}
\setlength{\parskip}{0pt}
\setlength{\headheight}{14pt}
\setlength{\footskip}{18pt}

\pagestyle{fancy}
\fancyhf{}
\fancyfoot[C]{\thepage}
\renewcommand{\headrulewidth}{0pt}
\captionsetup{font=normalsize,labelfont=bf,labelsep=period}
\renewcommand{\tablename}{Bảng}
\renewcommand{\refname}{Tài liệu tham khảo}
\urlstyle{same}

\titleformat*{\section}{\Large\bfseries}
\titleformat*{\subsection}{\large\bfseries}
\titleformat*{\subsubsection}{\normalsize\bfseries}
\newcolumntype{L}[1]{>{\raggedright\arraybackslash}p{#1}}
\newcolumntype{Y}{>{\raggedright\arraybackslash}X}

\begin{document}

\begin{center}
    	extbf{\Large BÁO CÁO ĐỀ XUẤT CẤU TRÚC VÀ CÔNG THỨC ĐO DOMAIN SHIFT}\\[0.2cm]
    	extbf{\Large CHO THÍ NGHIỆM NAS TRONG UNSUPERVISED ANOMALY DETECTION}\\[0.5cm]
    	extbf{Mục tiêu: viết lại báo cáo theo khung Vấn đề, Mục đích, Related Work, Method, Công thức; đồng thời thay representation shift chính từ MMD sang một lựa chọn mạnh hơn có reference A/A*}
\end{center}

\section{Vấn đề}

Trong thí nghiệm của đề tài, mô hình được huấn luyện trên dữ liệu nguồn và triển khai sang dữ liệu đích không có nhãn. Source và target thường vẫn thuộc cùng một họ hệ thống, nhưng khác nhau theo máy, điều kiện vận hành, môi trường, thời điểm đo, hiệu chuẩn, hoặc lão hóa. Vì vậy, bài toán cốt lõi không chỉ là ``mô hình có chạy được trên target hay không'', mà là phải lượng hóa được source--target mismatch theo cách đủ mạnh để giải thích khi nào adaptation thực sự có ích.

Về mặt xác suất, đây chủ yếu là bài toán \textbf{covariate shift dưới temporal nonstationarity}:
\[
P_s(X) \neq P_t(X).
\]

Concept drift
\[
P_s(Y \mid X) \neq P_t(Y \mid X)
\]
có thể tồn tại, nhưng không nên mặc định giả định nếu chưa có bằng chứng rằng ngữ nghĩa của normal/anomalous đã thay đổi. Với multivariate time series anomaly detection, dùng summary statistics thô hoặc domain classifier đặt trên mean--std features là chưa đủ, vì cách đó bỏ qua temporal context, latent geometry, và cấu trúc phụ thuộc giữa các sensor.

\section{Mục đích}

Bản báo cáo này theo đuổi bốn mục đích cụ thể:
\begin{enumerate}[label=\arabic*.]
    \item mô tả đúng bản chất domain shift trong setting source--target unlabeled deployment;
    \item chọn các reference chính chỉ từ các venue mạnh thuộc nhóm A/A* như AAAI, ICLR, ICML, NeurIPS, KDD;
    \item trả lời trực tiếp câu hỏi liệu có cách đo \textbf{representation shift} tốt hơn plain MMD hay không;
    \item đề xuất một protocol đo shift đủ chặt để hỗ trợ claim khi nào \texttt{adaptnas\_combined} có thể tốt hơn \texttt{uad\_source}.
\end{enumerate}

Kết luận ngắn gọn của báo cáo là: \textbf{có}. Nếu cần một representation-shift metric mạnh hơn MMD một lớp, nên dùng \textbf{JMMD} trên contextual features nhiều tầng theo tinh thần của JAN \cite{jan}. CORAL \cite{coral} giữ vai trò đo dependency shift bổ sung, còn DANN \cite{dann} chỉ nên dùng như sanity check về domain separability chứ không nên là metric chính.

\section{Related Work}

\subsection{Các paper nên giữ làm citation chính}

\small
\begin{longtable}{L{0.23\textwidth}L{0.13\textwidth}L{0.25\textwidth}L{0.31\textwidth}}
	oprule
	extbf{Bài báo} & \textbf{Venue} & \textbf{Nội dung chính} & \textbf{Kết luận dùng cho báo cáo} \\
\midrule
\endfirsthead
	oprule
	extbf{Bài báo} & \textbf{Venue} & \textbf{Nội dung chính} & \textbf{Kết luận dùng cho báo cáo} \\
\midrule
\endhead
OmniAnomaly \cite{omnianomaly} & KDD 2019 & Robust TSAD cho multivariate time series; SMD là benchmark quan trọng. & Hợp thức hóa bối cảnh TSAD trên SMD và việc đánh giá transfer giữa các machine là có ý nghĩa. \\
TS2Vec \cite{ts2vec} & AAAI 2022 & Học contextual representation cho time series ở mức timestamp và subsequence. & Domain shift nên được đo trong latent time-series space, không nên đo trên raw summary features. \\
CLUDA \cite{cluda} & ICLR 2023 & Học domain-invariant contextual representations cho UDA trên time series. & Củng cố luận điểm rằng contextual representation là đối tượng đúng để đo source--target mismatch. \\
JAN \cite{jan} & ICML 2017 & Căn chỉnh \emph{joint distributions} của nhiều tầng domain-specific bằng JMMD. & Nếu cần metric tốt hơn plain MMD, JMMD là lựa chọn mạnh hơn vì không chỉ nhìn một marginal distribution. \\
CORAL \cite{coral} & AAAI 2016 & Căn chỉnh second-order statistics giữa source và target. & Dependency shift nên được đo qua covariance mismatch trong latent space. \\
DANN \cite{dann} & ICML 2015 & Học feature vừa phân biệt nhiệm vụ vừa bất biến theo domain qua gradient reversal. & Domain separability nên là sanity check; nó hữu ích nhưng không đủ ổn định để làm metric chính. \\
\bottomrule
\end{longtable}
\normalsize

\subsection{Vì sao plain MMD không còn là lựa chọn chính}

Plain MMD vẫn là một baseline hợp lý về mặt kỹ thuật, nhưng trong báo cáo này nó không nên là trọng tâm vì hai lý do. Thứ nhất, reference nền tảng mạnh nhất của MMD là một bài journal, không khớp tiêu chí A/A* mà báo cáo đang muốn giữ thật sạch. Thứ hai, MMD một lớp chỉ so khớp \textbf{một marginal distribution} trong một latent space duy nhất; với multivariate time series, shift thường xuất hiện đồng thời ở nhiều mức ngữ cảnh và nhiều tầng biểu diễn. JAN \cite{jan} mạnh hơn ở điểm này vì nó so khớp \textbf{joint distributions của nhiều tầng}, nên phù hợp hơn để làm representation-shift metric chính.

\subsection{Lập trường chốt cho báo cáo}

Từ các paper trên, hướng viết chặt nhất là:
\begin{enumerate}[label=\arabic*.]
    \item dùng TS2Vec hoặc một encoder cùng tinh thần để tạo contextual embeddings;
    \item dùng JMMD làm \textbf{representation shift chính};
    \item dùng CORAL làm \textbf{dependency shift bổ sung};
    \item dùng DANN-style domain classifier chỉ làm \textbf{sanity check} về separability;
    \item nếu implementation hiện tại mới có một embedding cuối, có thể báo cáo MMD như baseline phụ, nhưng không nên để nó là công thức chính của phần Method.
\end{enumerate}

\section{Method}

\subsection{Dữ liệu dùng để đo shift}

Chỉ nên đo shift trên phần hành vi normal càng sạch càng tốt. Gọi \(X_s^N\) là tập source-normal windows, \(X_t^U\) là tập target unlabeled windows, và \(\widehat{X}_t^N \subseteq X_t^U\) là tập target windows được giữ lại sau một bước lọc bảo thủ nhằm giảm anomaly contamination. Mục đích của bước này là tránh để các anomaly hiếm nhưng mạnh chi phối shift score.

\subsection{Nguyên tắc chọn không gian biểu diễn}

Thay vì đo trên raw signal hoặc summary statistics, báo cáo này giả định có một encoder contextual \(f\) được huấn luyện theo tinh thần TS2Vec \cite{ts2vec} hoặc CLUDA \cite{cluda}. Nếu encoder có nhiều tầng hoặc nhiều mức ngữ cảnh, ta lấy một tập tầng \(\mathcal{L}\) để xây dựng representation shift. Đây chính là điểm làm JMMD phù hợp hơn MMD: shift không còn bị nén vào một embedding cuối duy nhất.

\subsection{Protocol đo shift được đề xuất}

Protocol đề xuất gồm ba thành phần:
\begin{enumerate}[label=\arabic*.]
    \item \textbf{Representation shift chính:} đo bằng JMMD trên các contextual features nhiều tầng.
    \item \textbf{Dependency shift bổ sung:} đo bằng CORAL distance trên embedding cuối hoặc embedding đã pooling.
    \item \textbf{Domain separability sanity check:} dùng domain-adversarial objective kiểu DANN để xác nhận source và target có thực sự dễ tách trong latent space hay không.
\end{enumerate}

Khuyến nghị báo cáo chính là giữ riêng hai score đầu tiên. Nếu cần xếp hạng source--target pairs, có thể chuẩn hóa rồi gộp ở phần thực nghiệm, nhưng phần Method nên nhấn mạnh rằng representation shift và dependency shift phản ánh hai khía cạnh khác nhau của mismatch, không nên collapse quá sớm thành một con số duy nhất.

\subsection{Trả lời trực tiếp câu hỏi ``có cách nào tốt hơn representation shift hiện tại không?''}

Câu trả lời là \textbf{có}. Nếu representation shift hiện tại đang được tính bằng MMD trên một embedding cuối, thì lựa chọn tốt hơn cho báo cáo là \textbf{JMMD trên nhiều tầng contextual features}. Lý do là:
\begin{enumerate}[label=\arabic*.]
    \item JMMD bắt được mismatch đồng thời ở nhiều tầng biểu diễn, thay vì chỉ một tầng;
    \item JMMD bám sát hơn với lập luận của deep transfer learning ở venue A* là ICML;
    \item JMMD phù hợp với multivariate time series, nơi temporal context và hierarchy của feature extractor đều mang thông tin về domain shift;
    \item CORAL vẫn nên đi kèm để bắt second-order dependency mismatch mà JMMD có thể không diễn giải trực tiếp.
\end{enumerate}

\section{Công thức, nguồn gốc công thức, nội dung paper, và kết luận}

\subsection{Công thức 1: Representation shift chính bằng JMMD}

\[
D_{\mathrm{rep}}^{\mathrm{JMMD}}(s,t)
= \left\|
\frac{1}{n_s}\sum_{i=1}^{n_s} \bigotimes_{\ell \in \mathcal{L}} \phi^{\ell}(h_{i,s}^{\ell})
-
\frac{1}{n_t}\sum_{j=1}^{n_t} \bigotimes_{\ell \in \mathcal{L}} \phi^{\ell}(h_{j,t}^{\ell})
\right\|^2.
\]

Trong đó, \(h_{i,s}^{\ell}\) và \(h_{j,t}^{\ell}\) là đặc trưng của source và target ở tầng \(\ell\), còn \(\phi^{\ell}(\cdot)\) là ánh xạ kernel tương ứng.

	extbf{Reference của công thức.} Long et al., \emph{Deep Transfer Learning with Joint Adaptation Networks}, ICML 2017 \cite{jan}.

	extbf{Paper nói gì.} JAN cho rằng nếu chỉ align một marginal distribution ở một layer cuối thì domain gap vẫn còn tồn tại trong các tầng task-specific. Vì vậy, paper đề xuất \textbf{Joint Maximum Mean Discrepancy} để align \textbf{joint distributions} của nhiều tầng cùng lúc.

	extbf{Kết luận của paper.} Trên benchmark UDA, joint adaptation tốt hơn các cách chỉ align một feature space đơn lẻ. Kết luận này rất quan trọng cho báo cáo hiện tại: nếu mục tiêu là đo representation shift nghiêm túc hơn plain MMD, thì \textbf{JMMD là lựa chọn tốt hơn MMD một lớp}.

	extbf{Kết luận rút ra cho bài toán này.} Với multivariate time series anomaly detection, shift không chỉ nằm ở embedding cuối mà còn nằm ở cách temporal context được mã hóa qua nhiều tầng. Vì vậy, báo cáo nên chốt \(D_{\mathrm{rep}}\) bằng JMMD là chính.

\subsection{Công thức 2: Dependency shift bằng CORAL distance}

\[
C_s = \mathrm{Cov}(Z_s),
\qquad
C_t = \mathrm{Cov}(Z_t),
\]
\[
D_{\mathrm{dep}}(s,t) = \frac{1}{4d^2}\left\| C_s - C_t \right\|_F^2.
\]

Trong đó, \(Z_s\) và \(Z_t\) là embedding cuối của source và target, còn \(d\) là số chiều embedding.

	extbf{Reference của công thức.} Sun et al., \emph{Return of Frustratingly Easy Domain Adaptation}, AAAI 2016 \cite{coral}.

	extbf{Paper nói gì.} CORAL lập luận rằng một phần quan trọng của domain shift nằm ở \textbf{second-order statistics}. Nếu covariance structure giữa hai domain khác nhau, mô hình huấn luyện ở source rất dễ suy giảm khi triển khai sang target. Paper vì vậy dùng chênh lệch covariance như một mục tiêu alignment rất đơn giản nhưng hiệu quả.

	extbf{Kết luận của paper.} Chỉ với việc align second-order statistics, một phương pháp rất đơn giản vẫn đạt kết quả cạnh tranh trên domain adaptation benchmark. Điều này cho thấy covariance mismatch là tín hiệu domain shift có giá trị thực sự, không phải chi tiết phụ.

	extbf{Kết luận rút ra cho bài toán này.} CORAL không thay thế representation shift, nhưng nó là metric bổ sung rất tốt để đo \textbf{dependency shift}. Trong TSAD đa biến, đây là phần cần giữ vì mismatch giữa source và target thường nằm ở coupling giữa các sensor, không chỉ ở mean hoặc variance riêng lẻ.

\subsection{Công thức 3: Domain-adversarial objective dùng làm sanity check}

\[
E(\theta_f,\theta_y,\theta_d)
= \mathcal{L}_y(\theta_f,\theta_y)
- \lambda \mathcal{L}_d(\theta_f,\theta_d).
\]

Trong đó, \(\mathcal{L}_y\) là loss của nhiệm vụ chính, \(\mathcal{L}_d\) là domain-classification loss, còn \(\theta_f, \theta_y, \theta_d\) lần lượt là tham số của feature extractor, label predictor, và domain discriminator.

	extbf{Reference của công thức.} Ganin and Lempitsky, \emph{Unsupervised Domain Adaptation by Backpropagation}, ICML 2015 \cite{dann}.

	extbf{Paper nói gì.} DANN tối ưu một feature extractor sao cho nó vừa giữ được tín hiệu phục vụ nhiệm vụ chính, vừa làm source và target khó phân biệt hơn đối với domain discriminator. Công cụ trung tâm là gradient reversal và objective đối kháng giữa tính phân biệt nhiệm vụ với tính bất biến theo domain.

	extbf{Kết luận của paper.} Học domain-invariant features bằng objective đối kháng cho kết quả adaptation mạnh trên các benchmark lớn. Tuy nhiên, điểm mạnh của DANN là \textbf{học representation thích nghi}, không phải cung cấp một metric ổn định để xếp hạng shift giữa mọi cặp domain.

	extbf{Kết luận rút ra cho bài toán này.} Domain classifier hữu ích để kiểm tra source và target có tách được trong latent space hay không, nhưng do nó phụ thuộc vào classifier capacity, tối ưu và regularization, nó chỉ nên là \textbf{sanity check}, không nên là công thức chính của shift score.

\section{Kết luận chốt để dùng trong báo cáo}

Nếu cần viết ngắn gọn, chặt và đúng tiêu chí A/A*, phần Method nên chốt như sau:
\begin{enumerate}[label=\arabic*.]
    \item bài toán được xem là covariate shift dưới temporal nonstationarity;
    \item không đo shift trên raw summary features;
    \item dùng contextual representation theo tinh thần TS2Vec \cite{ts2vec} và CLUDA \cite{cluda};
    \item dùng \textbf{JMMD} \cite{jan} làm representation shift chính vì nó mạnh hơn plain MMD một lớp;
    \item dùng \textbf{CORAL} \cite{coral} làm dependency shift bổ sung;
    \item dùng \textbf{DANN-style domain separability} \cite{dann} chỉ như sanity check;
    \item khi trình bày kết quả, nên ưu tiên báo cáo riêng \(D_{\mathrm{rep}}^{\mathrm{JMMD}}\) và \(D_{\mathrm{dep}}\), sau đó mới xét việc gộp score nếu cần ranking.
\end{enumerate}

Tóm lại, câu trả lời cho câu hỏi mở đầu là: \textbf{có}. Nếu muốn tính representation shift tốt hơn, báo cáo nên chuyển từ \textbf{MMD trên một embedding cuối} sang \textbf{JMMD trên nhiều tầng contextual features}; đây là lựa chọn vừa mạnh hơn về mặt biểu diễn, vừa sạch hơn về mặt citation theo tiêu chí A/A*.

\begin{thebibliography}{99}

\bibitem{omnianomaly}
Yonghong Su, Yuliang Zhao, Chongzhou Niu, Rong Liu, Wei Sun, and Dan Pei.
\newblock OmniAnomaly: Robust Anomaly Detection for Multivariate Time Series through Stochastic Recurrent Neural Network.
\newblock In \emph{Proceedings of the 25th ACM SIGKDD Conference on Knowledge Discovery and Data Mining}, 2019.

\bibitem{ts2vec}
Zhihan Yue, Yujing Wang, Juanyong Duan, Tianmeng Yang, Congrui Huang, Yunhai Tong, and Bixiong Xu.
\newblock TS2Vec: Towards Universal Representation of Time Series.
\newblock In \emph{Proceedings of the AAAI Conference on Artificial Intelligence}, 2022.

\bibitem{cluda}
Yilmazcan Ozyurt, Stefan Feuerriegel, and Ce Zhang.
\newblock Contrastive Learning for Unsupervised Domain Adaptation of Time Series.
\newblock In \emph{International Conference on Learning Representations}, 2023.

\bibitem{dann}
Yaroslav Ganin and Victor Lempitsky.
\newblock Unsupervised Domain Adaptation by Backpropagation.
\newblock In \emph{Proceedings of the 32nd International Conference on Machine Learning}, 2015.

\bibitem{jan}
Mingsheng Long, Han Zhu, Jianmin Wang, and Michael I. Jordan.
\newblock Deep Transfer Learning with Joint Adaptation Networks.
\newblock In \emph{Proceedings of the 34th International Conference on Machine Learning}, 2017.

\bibitem{coral}
Baochen Sun, Jiashi Feng, and Kate Saenko.
\newblock Return of Frustratingly Easy Domain Adaptation.
\newblock In \emph{Proceedings of the AAAI Conference on Artificial Intelligence}, 2016.

\end{thebibliography}

\end{document}
